\documentclass{article}

\usepackage{nips_2018_author_response}

\usepackage[utf8]{inputenc} % allow utf-8 input
\usepackage[T1]{fontenc}    % use 8-bit T1 fonts
\usepackage{hyperref}       % hyperlinks
\usepackage{url}            % simple URL typesetting
\usepackage{booktabs}       % professional-quality tables
\usepackage{amsfonts}       % blackboard math symbols
\usepackage{nicefrac}       % compact symbols for 1/2, etc.
\usepackage{microtype}      % microtypography

\usepackage[rightcaption]{sidecap}
\usepackage{graphicx}
\usepackage{verbatim}
\begin{document}

We thank all of our reviewers for their thoughtful and detailed comments, and we are revising our our manuscript in light of this helpful feedback.
Here, we discuss a handful of specific points:
%\subsubsection*{Choice of initial programming primitives}

\textbf{Choice of initial programming primitives:} Reviewer 1 observes that our initial set of programming primitives is
not as minimal as it could have been,
raising intriguing questions
about the importance and role of the initial basis.
We believe that general program-learners should not start with a \emph{minimal} basis,
but rather a \emph{domain-agnostic}
basis,
like those embodied in the standard libraries of modern programming languages.
An open question however is the extent to which one could ``learn from scratch.''
As a step toward answering this question,
we ran the following experiment:
SCC was provided with a minimal subset of primitives in 1959 Lisp, crucially including primitive recursion,
and tasked with solving a moderate set of functional programming exercises, like appending lists or summing the elements of a list.
After 93 hours on 64 CPUs,
SCC rediscovered \texttt{fold}, \texttt{unfold}, \texttt{map},
\texttt{length}, \texttt{zip}, \texttt{range}, \texttt{index},
and a handful of arithmetic operations.
(We also tried the exact experiment Reviewer 1 suggests---removing only \texttt{length}---and confirmed that SCC rediscovers \texttt{length} when trained on our list processing domain.)
Although this result is suggestive,
we believe that in practice learners should start from a rich yet domain-agnostic
inventory of primitives, like human programmers do.


%% As several reviewers pointed out the original \texttt{DSL} influences the efficiency and
%% accuracy of the algorithm. Reviewer 1 points out ``What happens, for example, if
%% length is removed [from the original list \texttt{DSL}]?''. To explore that question we
%% took the original \texttt{Lisp} specification from [citation needed] as well as
%% tasks from \texttt{Lisp} textbooks to see whether higher level functions would
%% be discovered. While this example was computationally expensive, it successfully
%% rediscovered common abstractions from the functional programming world,
%% including but not limited to \texttt{fold}, \texttt{unfold}, \texttt{map},
%% \texttt{length}, \texttt{zip}, \texttt{range}, etc.

%% That being said, the intended goal of the algorithm is not to rediscover
%% well-known useful abstractions but rather to adapt such abstractions to the
%% tasks at hands, through more abstractions if need be but not with that as a goal
%% \textit{per se}.

\textbf{The evolution of task solutions over time:} %\subsubsection*{The evolution of task solutions over time}
Reviewer 1 asks about the ``evolution of solutions to a fixed task through the iterations of the main algorithm.''
In most cases the solution to a fixed task
stays the same,
modulo rewriting the solution in terms of new subroutines.
There are interesting cases however where the
semantics of the program changes as the algorithms constructs more subroutines. Figure~1
illustrates one case for turtle graphics.
\begin{SCfigure}[2][h]
  \centering
  \includegraphics[width=7cm]{figures/rebutal/spirals.eps}
  \caption{Three turtle graphics tasks: the task is to write a program that draws the image. Initially SCC solves the leftmost task without using a loop, but after solving the middle and rightmost tasks,
  it constructs a parametric `spiral' subroutine, which uses a loop -- on subsequent iterations it solves the leftmost task using the `spiral' routine.}
\end{SCfigure}

\textbf{Fragment Grammars:} Reviewers 1 \& 2 both request clarification on the mechanics of the DSL induction,
which employs a representation from computational linguistics called Fragment Grammars (O'Donnell 2015).
We are revising the Supplement to include a brief overview of this formalism.

\textbf{Reproducibility:} %\subsubsection*{Reproducibility}
Our source code is publicly available on GitHub,
and a link will be added upon publication (in NIPS or elsewhere),
along with the specific command lines  needed to run
the experiments in the paper.

\textbf{Specific questions from Reviewer 2:}
The 236 list tasks were created by one of the authors of this paper
several months before this project began
and without this system (SCC) in mind,
with the goal of
being human-interpretable
yet reasonably challenging for
program synthesizers.
The AIC (Akaike Information Criterion) is a model comparison criteria 
similar to the BIC (Bishop 2006: ``Pattern Recognition and Machine Learning'' has an overview),
which introduces the counting norm term in Eq.~3.
The $\theta_{var}$ parameter
is the probability that
a program sampled from the prior will
have a variable at a particular node it syntax tree -- this is specified in Algorithm 1 of the supplement,
but we admit that this is quite buried,
so we will include a pointer to this part of the supplement in the revised manuscript.
We compare our full system against an ablation 
without the recognition model,
finding that
the full system solves more programming tasks and takes less time to
solve those tasks.

\textbf{FlashFill:} Reviewer 2 suggests comparing with FlashFill, the state-of-the-art program synthesizer for text editing.
We very much would like to do this comparison, but this comparison
is encumbered by the fact that FlashFill is not freely available nor accessible outside Microsoft Excel.
We strongly believe that FlashFill would easily solve essentially all of our text editing problems.
Instead, we compared with the winner of the 2017 SyGuS program synthesis competition (Alur et al. 2017), 
showing that, without hand-engineering of the DSL,
we achieve results comparable to the 2017 winning solver (CVC4),
whereas CVC4 was provided with hand-engineered DSLs (Section 3.2 of paper).
%\subsubsection*{Related Work}

\textbf{Related Work:} We thank reviewer three for suggesting points of contact between our algorithm and other work in program induction.
We are revising our paper to prominently cite the work of Schmid and Kitzelmann,
whose \textsc{Igor} systems, like ours,
are motivated by the flexibility and generality of human learning.
Katayama's \textsc{MagicHaskeller} synthesizer for functional programs will also be cited.
More broadly,
we see these  related works as
 complementary to
our system:
we take a black-box synthesizer and
learn to use it by growing the DSL and
training a recognition model,
and one could replace our enumerative synthesizer with
\textsc{Igor}, Sketch (Solar-Lezama 2008), \textsc{MagicHaskeller},
genetic programming, or other synthesis algorithms.
We chose however to
use a generic and simple enumerative approach because it allowed us to apply SCC
to a diverse spread of problems: functional programming exercises, text editing,
turtle graphics programming, and symbolic regression from images.

%% revision will prominently highlight the work of Kitzelmann, Schmid and Katayama,
%% particularly Kitzelmann and Schmid, whose cognitive inspiration
%% is closely allied with our motivations.

%% Kitzelmann and Schmid's \textsc{Igor} systems are sophisticated
%% analytical inductive programming algorithms that handle certain classes of
%% problems, much like {Sketch} and {FlashFill} are sophisticated program
%% synthesizers for certain classes of programs.
%% We see these lines of work as orthogonal and complementary to that of SCC\@. Our
%% goal is to take a black-box synthesizer and then learn how to use it to solve a class of
%% problems. One could replace our enumerative synthesizer
%% with \textsc{Igor}, Sketch or \textsc{MagicHaskeller}. 

%% This also relates to an insightful
%% and suggestion for comparison by reviewer 2: FlashFill's DSL is fixed and
%% manually curated and enhanced for a specific class of problems, over which is
%% performs well at the loss of agility.


%j\subsubsection*{More specific questions}

%% \paragraph*{Reviewer 2: Origin of list task} All the list tasks were manually
%% created in order to be interesting, challeging, but also human-interpretable for
%% a prior behaviour experiment. They pre-existed our algorithm.



%% Our project --- including source code, data and automation scripts --- is
%% already available on \url{GitHub} and future revision will include a link to it
%% once the anonymity of the authors can be unveiled. No technical knowledge will
%% be required to reproduced the results.

\begin{comment}
  
Reviewer #1
Questions

    1. Please provide an "overall score" for this submission.
        8: Top 50% of accepted NIPS papers. A very good submission; a clear accept. I vote and argue for accepting this submission.
    2. Please provide a "confidence score" for your assessment of this submission.
        4: You are confident in your assessment, but not absolutely certain. It is unlikely, but not impossible, that you did not understand some parts of the submission or that you are unfamiliar with some pieces of related work.
    3. Please provide detailed comments that explain your "overall score" and "confidence score" for this submission. You should summarize the main ideas of the submission and relate these ideas to previous work at NIPS and in other archival conferences and journals. You should then summarize the strengths and weaknesses of the submission, focusing on each of the following four criteria: quality, clarity, originality, and significance.
        = Summary

        A method for learning a DSL for program synthesis together with a search
        algorithm in that DSL is presented. The method proceeds iteratively,
        trying to solve tasks with the current DSL, and then extracting new DSL
        components from the solutions. Experiments show that bootstrapping the
        method with a DSL made up of trivial primitives is sufficient to
        discover common high-level constructs present in manually constructed
        DSLs.

        The paper tackles an important problem (DSL design) in an elegant and
        novel way. The clarity of the paper is not perfect, as the details of the
        idea require more space than the 8 pages available, but it clearly is
        stepping stone towards a new generation of program synthesis approaches.

        = Quality

        The paper precisely defines its main algorithm and the core concepts of
        the approach. There are a few omissions (e.g., Sect. 2.4 does not
        explain the exact mechanism to extract new constructs; the supplement
        gives an algorithm, but omits a clear definition of how "fragments" are
        found), but overall the paper is precise and explains core concepts well.
        The experiments are very thorough, comparing the method to interesting
        ablations and baselines.
        What I am missing is an analysis of the influence of the initial set of
        primitives. For example, the list processing set starts with a
        non-minimal set of primitives (as length is "foldr (\lambda (n) -> 1 + n) 0"
        and index can be implemented using the others as well, though
        painfully). What happens, for example, if length is removed?

        As an additional data point, I would have appreciated a look at the
        "evolution" of solutions to a fixed task through the iterations of the
        main algorithm, i.e., show an initial, early and late solution in the
        style of Table 1.

        = Clarity

        Many part of the paper are better understood with the supplement in hand
        (Sect. 2.3 is not discussing details of how q is modeled, but Supplement
        Sect. 5 makes clear that is is essentially like DeepCoder; Sect. 2.4 is
        missing details as discussed above).
        Overall, the paper is easily understood, but not entirely self-contained
        and not precise enough to allow re-implementation by third parties.

        The supplement is very unpolished (there are many grammatical errors and
        half-finished sentences, e.g. in line 6, line 33, lines 45/46, ...)

        = Originality

        While drawing on the program synthesis literature of the last few years
        (e.g., the guided enumerative search), the core contributions of the
        paper are completely new and are a substantial step forward for the
        field.

        = Significance

        Designing a suitable DSL is often one of the hardest tasks in applying
        program synthesis to a new domain; the promise of automatically learning
        a DSL (together with a search algorithm) is obvious. I am not entirely
        convinced by the authors idea of generalizing to other generative
        programs, but I am looking forward to being proven wrong.
    4. How confident are you that this submission could be reproduced by others, assuming equal access to data and resources?
        2: Somewhat confident

Reviewer #2
Questions

    1. Please provide an "overall score" for this submission.
        8: Top 50% of accepted NIPS papers. A very good submission; a clear accept. I vote and argue for accepting this submission.
    2. Please provide a "confidence score" for your assessment of this submission.
        3: You are fairly confident in your assessment. It is possible that you did not understand some parts of the submission or that you are unfamiliar with some pieces of related work. Math/other details were not carefully checked.
    3. Please provide detailed comments that explain your "overall score" and "confidence score" for this submission. You should summarize the main ideas of the submission and relate these ideas to previous work at NIPS and in other archival conferences and journals. You should then summarize the strengths and weaknesses of the submission, focusing on each of the following four criteria: quality, clarity, originality, and significance.
        This paper proposes an algorithm for program synthesis in an incremental way. There are three steps to it: 1) searching for program solutions, 2) expandind the dsl by discovering useful reusable components, 3) training a neural network that primes the program solution search in 1). The interesting part of the contribution is the focus on roughly mimicking the way humans refactor, recompose and reuse when coding.

        Strengths
        - the problem presented is interesting, and the scc algorithm is novel. the theoretical framework is sound, and approximations made are reasonable and useful
        - the paper presents a good evaluation framework on three different tasks and through the ablation analysis that clearly shows the contributions of each part of the model to the final results
        - the results are impressive - excellent performance on all three tasks, both accuracy-wise and time-wise


        Weaknesses
        - this is quite a dense paper and it packs a lot of information
        - the paper could benefit from a comparison to a well known systems such as flashfill, where the dsl is fixed

        Questions:
        - 125-126 the probability of a variable occurring in a program. could you provide an example here?
        - equation 3 - is this the counting norm? if so, can you provide a quick rationale why it is used
        - 167 - can you explicitly say what AIC is and/or cite it?
        - 181-183 - how is this unification done?
        - 200 - was this manually created?
        - what was the recognition model validated on?

        Other:
        - figure 4 position is weird
        - figure 4 bottom left and right are not clear
    4. How confident are you that this submission could be reproduced by others, assuming equal access to data and resources?
        1: Not confident

Reviewer #3
Questions

    1. Please provide an "overall score" for this submission.
        7: A good submission; an accept. I vote for accepting this submission, although I would not be upset if it were rejected.
    2. Please provide a "confidence score" for your assessment of this submission.
        3: You are fairly confident in your assessment. It is possible that you did not understand some parts of the submission or that you are unfamiliar with some pieces of related work. Math/other details were not carefully checked.
    3. Please provide detailed comments that explain your "overall score" and "confidence score" for this submission. You should summarize the main ideas of the submission and relate these ideas to previous work at NIPS and in other archival conferences and journals. You should then summarize the strengths and weaknesses of the submission, focusing on each of the following four criteria: quality, clarity, originality, and significance.
        The authors present an inductive programming approach where a DSL is learned by a search, compile, and compress strategy. They evaluate their approach in three different domains -- list processing, text editing, and symbolic regression. They motivate their approach by pointing out the effectivenes and flexibility of human learning. Similar arguments are presented by Schmid & Kitzelmann, Cognitive Systems Research 2011 and the approach has some similarities with the MagicHaskeler of Katayama.Although some relation to pevious research is missing (mainly the two groups named above), the presented approach is original. I am most impressed by the empirical demonstration of applicability in three rather distinct domains (but: see also Schmid/Kitzelmann, 2011). The presentation is technically sound and the empirical evaluation is very convincing. 
    4. How confident are you that this submission could be reproduced by others, assuming equal access to data and resources?
        3: Very confident


  \end{comment}


\end{document}
